\chapter{Simulación en Gazebo}\label{cap:sim}

El paquete \cod{simulacion} (\archivo{ros2\_ws/src/simulacion/}) arma en \textbf{Gazebo Harmonic} una hilera de
limoneros y el robot, con \textbf{los mismos tópicos y marcos que los ESP32}. Así \cod{mision}, \cod{brazo\_ik} y la
visión corren sin cambios, sin hardware.

\section{Qué hay en el mundo}
\begin{itemize}
  \item Pasillo a lo largo de $+x$, robot en el origen, \textbf{4 limoneros por lado} en $x$ = 1,5 / 3 / 4,5 / 6~m,
        $y = \pm 0{,}70$~m, con las medidas de R1.
  \item \textbf{12 limones cosechables} (3 por árbol de la fila izquierda).
  \item 2 obstáculos en el pasillo (una caja y un balde) que obligan a esquivar.
  \item \textbf{Marcas de parada} blancas en $x$ = 1,19 / 2,70 / 4,17 / 5,69~m: frenando con el eje de las ruedas
        sobre una marca, en el centro del pasillo, los 3 limones de ese árbol quedan a la vista y al alcance.
\end{itemize}

\section{Quién imita a quién}
{\small
\begin{tabularx}{\textwidth}{@{}l l X@{}}
\toprule
\textbf{Nodo de la simulación} & \textbf{Reemplaza a} & \textbf{Detalle} \\
\midrule
\cod{base\_sim} & firmware de la S3 & misma seguridad (\cod{SeguridadBase} portada), arranca DESARMADA \\
\cod{brazo\_sim} & firmware de la WROOM & mismos tópicos, perfil trapezoidal, arranca DESARMADO \\
\cod{limones\_sim} & visión & publica \topic{/limon/posicion} solo con el robot quieto y el limón en la zona útil \\
3 \cod{gpu\_lidar} & HC-SR04 & 5 rayos en un abanico de 15° cada uno \\
3 planchas con resorte & bumpers & apretado = hundida más de 3~mm \\
\bottomrule
\end{tabularx}}

\section{Cómo se usa}

\begin{lstlisting}[language=consola]
cd ~/Documentos/recolector_de_frutas/ros2_ws && source /opt/ros/jazzy/setup.bash
colcon build --symlink-install --packages-select simulacion && source install/setup.bash
export ROS_AUTOMATIC_DISCOVERY_RANGE=LOCALHOST          # ¡en CADA terminal! (aislar de los ESP32 reales)

ros2 launch simulacion sim.launch.py                    # robot y hilera (gui:=false: sin ventana)
ros2 launch simulacion sim.launch.py mision:=true       # + misión (arrimar con el teclado a una marca)
ros2 launch simulacion sim.launch.py mision:=true autonomo:=true   # cosecha sola toda la hilera

# Manejar a mano (otra terminal, con el export):
ros2 topic pub --once /base/comando std_msgs/msg/String "{data: armar}"
ros2 run teleop_twist_keyboard teleop_twist_keyboard --ros-args -p repeat_rate:=10.0
\end{lstlisting}

\begin{nota}[\cod{repeat\_rate}]
Sin \cod{repeat\_rate:=10.0}, el teclado publica solo al tocar una tecla y el watchdog de 300~ms frena la base entre
tecla y tecla.
\end{nota}

\begin{peligro}[Aislar la simulación]
Sin \cod{ROS\_AUTOMATIC\_DISCOVERY\_RANGE=LOCALHOST}, la simulación ve por la red los ESP32 reales conectados a la Pi:
el 28/09 la misión simulada leyó el estado del brazo real y le mandó comandos.
\end{peligro}

Si algo queda colgado después de cortar:
\begin{lstlisting}[language=consola]
pkill -9 -f "[g]z sim"
pkill -9 -f "[b]razo_sim|[b]ase_sim|[l]imones_sim|[p]arameter_bridge"
\end{lstlisting}

\section{Resultados}
\begin{itemize}
  \item Con \cod{mision:=true autonomo:=true}: recorre la hilera hasta $x = 7{,}32$~m en $\sim$9~min simulados y
        cosecha \textbf{12 de 12} limones, esquivando la caja por la derecha y el balde por la izquierda, sin tocarlos.
  \item El watchdog frena la base a los 300~ms (11--13~cm desde 0,2~m/s con la rampa).
  \item Con la cadena de visión real (\cod{vision:=yolo}): el estéreo localiza, pero YOLO no reconoce los limones
        sintéticos. No reentrenar con imágenes de Gazebo: arruinaría el modelo para los limones reales.
\end{itemize}

\begin{nota}[Trucos de la física]
Las masas del brazo van al 5~\% (si no, los impulsos del brazo deslizaban la base con las ruedas frenadas), el limón
agarrado se lleva con \cod{set\_pose} y \cod{limones\_sim} usa la pose real de Gazebo, no la odometría. Son cosas de la
simulación: en el robot real no aplican.
\end{nota}
